\chapter{Wnioski}
\label{chap:wnioski}

\section{Podsumowanie realizacji pracy}

Celem pracy było empiryczne zbadanie granic wydajności transmisji bezprzewodowej w~systemach wbudowanych czasu rzeczywistego oraz wskazanie kompromisów, które należy uwzględnić przy wyborze protokołu komunikacyjnego.
Cel ten zrealizowano poprzez zaprojektowanie, implementację i~wykorzystanie autorskiego środowiska badawczego ESPy-Nexus, działającego w~architekturze \textit{Hardware-in-the-Loop}.

Opracowane środowisko obejmuje trzy współpracujące warstwy.
Oprogramowanie wbudowane uruchamiane na mikrokontrolerze ESP32 realizuje obsługę płaszczyzny sterowania, akwizycję danych oraz rejestrację stempli czasowych.
Zastosowanie zewnętrznej pamięci PSRAM pozwoliło przy tym na buforowanie dużej liczby rekordów, natomiast binarna akwizycja ograniczyła narzut przesyłania wyników z~ESP32 do komputera.
Silnik testowy w~języku Python odpowiada za konfigurację eksperymentów, generowanie ruchu, komunikację z~urządzeniem i~zapis wyników w~bazie SQLite.
Uzupełnieniem systemu jest aplikacja webowa wykonana z~wykorzystaniem frameworka Nuxt, która umożliwia filtrowanie, porównywanie i~wizualizację zagregowanych metryk.

W~ramach pracy przygotowano macierz obejmującą topologie \texttt{NONE}, \texttt{AP\_ESP}, \texttt{AP\_PC} i~\texttt{AP\_ROUTER}, protokoły \texttt{SERIAL\_STR}, \texttt{SERIAL\_BIN}, \texttt{UDP}, \texttt{TCP} i~\texttt{WS}, ładunki o~rozmiarach od 4~B~do 1024~B~oraz częstotliwości od 1~Hz do 10.000~Hz.
Dane surowe były następnie przekształcane do ujednoliconych wskaźników jakości transmisji, takich jak PDR, jitter, czas przerwy w~odbiorze, goodput, liczba zdarzeń \textit{Out-of-Order} i~względne opóźnienie.

\section{Stopień osiągnięcia celu}

Założony cel pracy został osiągnięty w~wysokim stopniu.
Zbudowany system umożliwił przeprowadzenie powtarzalnych eksperymentów z~rzeczywistym mikrokontrolerem.
Dzięki temu w~wynikach uwzględniono ograniczenia procesora, pamięci, sterowników i~stosów komunikacyjnych, których nie można w~pełni odtworzyć w~symulacji programowej.

Zrealizowano również cel analityczny.
Zestaw metryk pozwolił rozdzielić trzy różne aspekty jakości transmisji: niezawodność dostarczania, stabilność czasową oraz przepustowość użyteczną.
Weryfikacja na scenariuszach mock potwierdziła poprawne rozpoznawanie strat, duplikatów, nieregularności odstępów między pakietami, zmian kolejności i~narastającego opóźnienia.
Jednocześnie analiza wykazała, że część wskaźników, zwłaszcza względne opóźnienie i~dryft wyrażony w~PPM, należy interpretować jako miary diagnostyczne zależne od przyjętego modelu, a~nie jako bezpośredni pomiar pojedynczego zjawiska fizycznego.

Trzeba jednak zastrzec, że nie można uznać systemu za deterministyczne źródło czasu rzeczywistego.
Generator ruchu działa na komputerze z~systemem GPOS, dlatego aktywne oczekiwanie ogranicza, lecz nie usuwa wpływu planisty zadań, sterowników i~buforów systemowych.
Ograniczenie to zostało zidentyfikowane i~uwzględnione w~interpretacji wyników, dlatego nie podważa realizacji celu projektowego, lecz wyznacza zakres wiarygodnych wniosków.

\section{Najważniejsze wyniki badań}

Poniżej zebrano najważniejsze wnioski, które wynikają wprost z~przeprowadzonej kampanii pomiarowej:

\begin{enumerate}
    \item Nie istnieje jeden protokół komunikacyjny, który byłby najlepszy dla wszystkich badanych kryteriów.
          Wybór rozwiązania powinien wynikać z~wymagań dotyczących dopuszczalnych strat, opóźnienia, zmienności czasowej, przepustowości oraz rozmiaru przesyłanych danych.

    \item Interfejs szeregowy może stanowić użyteczną linię bazową, lecz jego przepustowość szybko ogranicza wzrost rozmiaru ładunku.
          Dla \texttt{SERIAL\_BIN} z~ładunkiem 4~B~PDR utrzymywał się na poziomie 100\% aż do 5600~Hz, a~dopiero przy 5700~Hz spadł do 93,4\%.
          Dla ładunku 1024~B~pierwsze straty pojawiły się już przy 40~Hz, a~przy 50~Hz PDR wyniósł 1,1\%.

    \item TCP i~WebSockets skutecznie chronią integralność uporządkowanego strumienia, lecz wysoki PDR nie oznacza niskiego opóźnienia.
          W~badanym przypadku TCP zachowywał PDR równy 100\%, ale przy dużym obciążeniu występowały opóźnienia względne rzędu sekund oraz spadek goodputu.
          Dla konfiguracji TCP, \texttt{AP\_PC}, 4~B~i~3900~Hz odnotowano CV około 3000\% oraz maksymalne opóźnienie względne około 2,23~s.

    \item UDP oferuje korzystny kompromis dla zastosowań, w~których ważniejsze od dostarczenia każdego pakietu jest ograniczenie opóźnienia i~obsługa aktualnych danych.
          W~konfiguracji \texttt{AP\_ROUTER} i~dla ładunku 1024~B~goodput osiągnął około 28,03~Mbps, jednak przy tym obciążeniu PDR spadł do około 65\%.
          Wynik pokazuje, że wzrost przepustowości może odbywać się kosztem niezawodności.

    \item Topologia sieciowa istotnie wpływa na wyniki.
          Dla UDP i~ładunku 1024~B~pierwsze obserwowane straty wystąpiły przy około 100~Hz w~topologii \texttt{AP\_PC}, 128~Hz w~topologii \texttt{AP\_ROUTER} i~300~Hz w~topologii \texttt{AP\_ESP}.
          Przy wysokich częstotliwościach ranking mógł się jednak odwracać: dla ładunku 16~B~i~zakresu 9000-10.000~Hz średni PDR wyniósł około 91,6\% dla \texttt{AP\_PC} oraz 26,6\% dla \texttt{AP\_ESP}.

    \item Wskaźnik PDR nie jest wystarczający do oceny przydatności transmisji w~systemach czasu rzeczywistego.
          Kompletne dostarczenie danych może współwystępować z~opóźnieniem uniemożliwiającym terminową reakcję, a~utrata części pakietów może być akceptowalna, jeśli aplikacja toleruje brak pojedynczych aktualizacji i~wymaga małego opóźnienia.
\end{enumerate}

\section{Zalecenia projektowe}

Z~powyższych wyników wynikają następujące zalecenia projektowe:

\begin{itemize}
    \item W~przypadku zamkniętej pętli sterowania, w~której najważniejsze są aktualność danych i~ograniczenie opóźnienia, należy w~pierwszej kolejności rozważyć UDP.
          Protokół powinien zostać uzupełniony o~warstwę aplikacyjną odpowiadającą za identyfikację pakietów, kontrolę ich aktualności, odrzucanie przestarzałych ramek oraz reakcję na utratę danych.

    \item TCP lub WebSockets należy stosować wtedy, gdy kluczowe są integralność, uporządkowanie i~dostarczenie całego strumienia, na przykład podczas przesyłania konfiguracji, plików lub aktualizacji oprogramowania.
          Protokoły te nie powinny być wybierane do szybkiego sterowania bez dodatkowego mechanizmu kontroli maksymalnego opóźnienia.

    \item Interfejs UART należy traktować jako rozwiązanie dla krótkich komend, konfiguracji i~diagnostyki albo jako punkt odniesienia.
          Przy większych ładunkach trzeba uwzględnić ograniczoną przepływność, buforowanie oraz możliwość długotrwałej przerwy w~odbiorze.

    \item Dobór topologii powinien być wykonywany dla rzeczywistego zakresu częstotliwości i~rozmiaru ładunku.
          Wynik uzyskany dla małego obciążenia nie pozwala przewidzieć zachowania przy nasyceniu, dlatego rekomendowane jest badanie co najmniej kilku punktów poniżej i~powyżej oczekiwanego progu degradacji.

    \item Kryterium akceptacji transmisji powinno obejmować jednocześnie PDR, jitter, maksymalny czas przerwy w~odbiorze, goodput oraz względne opóźnienie.
          W~praktyce należy zdefiniować limity dla każdej z~tych metryk przed rozpoczęciem kampanii pomiarowej.

    \item Przy testowaniu z~użyciem komputera z~systemem GPOS należy rejestrować rzeczywisty czas nadawania i~traktować aktywne oczekiwanie jako metodę ograniczania jitteru, a~nie jako gwarancję determinizmu.
          Wyniki bliskie granicy akceptacji powinny być potwierdzane dodatkowymi powtórzeniami.
\end{itemize}

\section{Ograniczenia i~kierunki dalszych badań}

Najważniejszym ograniczeniem pracy jest wykorzystanie komputera z~systemem operacyjnym ogólnego przeznaczenia jako generatora ruchu.
Otrzymane wyniki opisują zatem cały badany tor transmisyjny, obejmujący komputer, sterowniki, interfejs sieciowy, urządzenia pośredniczące i~ESP32, a~nie wyłącznie właściwości samego protokołu.
Dodatkowo część wniosków opiera się na pojedynczej platformie sprzętowej, określonym firmware oraz konkretnej konfiguracji topologii.

Naturalnym rozwinięciem pracy jest więc rozszerzenie badań o~kolejne warianty sprzętowe i~metodologiczne.
W~dalszych badaniach zasadne jest:

\begin{itemize}
    \item powtórzenie kampanii na innych modułach ESP32, kartach sieciowych, routerach i~wersjach systemu operacyjnego;
    \item porównanie generatora działającego w~GPOS z~generatorem czasu rzeczywistego lub sprzętowym generatorem ramek;
    \item zastosowanie synchronizacji zegarów, na przykład przez PTP, w~celu dokładniejszego pomiaru opóźnienia jednokierunkowego;
    \item wydłużenie pojedynczych testów i~wielokrotne powtórzenie tych samych konfiguracji w~celu wyznaczenia przedziałów niepewności oraz rozrzutu wyników;
    \item uzupełnienie analizy o~temperaturę, zużycie pamięci, obciążenie procesora, stan buforów oraz informacje z~warstw niższych, aby lepiej rozdzielić źródła strat i~opóźnień;
    \item rozszerzenie systemu o~automatyczne wyznaczanie progów akceptacji i~raportowanie konfiguracji spełniających wymagania konkretnej aplikacji.
\end{itemize}

\section{Wniosek końcowy}

ESPy-Nexus spełnił założoną funkcję jako zautomatyzowane środowisko do badania rzeczywistego urządzenia ESP32 w~pętli sprzętowej.
Połączenie warstwy firmware, silnika testowego, analizy ETL i~interfejsu webowego umożliwiło przejście od generowania ruchu do porównywalnych wskaźników jakości transmisji.

Przeprowadzone badania potwierdziły, że ocena protokołu wyłącznie na podstawie PDR prowadzi do niepełnych wniosków.
UDP może zapewniać wysoką przepustowość i~ograniczone opóźnienia kosztem strat, natomiast TCP i~WebSockets mogą dostarczać komplet danych kosztem opóźnienia i~zmienności czasowej.
Ostateczny wybór powinien więc wynikać z~wymagań aplikacji oraz z~pomiarów wykonanych w~docelowej topologii, a~nie z~ogólnej klasyfikacji protokołów.
